\subsection{\texorpdfstring{EULERIAN EXPANSION OF \(\cot z\)}{EULERIAN EXPANSION OF COT Z}}

By formula (20) of VI, p.~275, the numbers \(b_{n} / n!\) are the coefficients in the expansion of \(S / (e^{S} - 1)\) as a formal series; we shall show in this section that the function \(z / (e^{z} - 1)\) is equal to the sum of an absolutely convergent entire series on a neighbourhood of 0 in \(\mathbf{C}\) ; it will follow from the lemma of VI, p.~280, that the coefficients of this series are the numbers \(b_{n} / n!\) , from which we shall deduce estimates for the Bernoulli numbers \(b_{n}\) .

In the first place we note that

\[
{\frac {z}{e ^ {z} - 1}} = - {\frac {z}{2}} + {\frac {z}{2}} {\frac {e ^ {z} + 1}{e ^ {z} - 1}} = - {\frac {z}{2}} + {\frac {i z}{2}} \cot {\frac {i z}{2}}.\tag{1}
\]

We shall obtain below a series expansion for \(\cot z\) , valid for every \(z\) not an integer multiple of \(\pi\) .

PROPOSITION 1. For every complex number \(z\) and every integer \(n\) one has

\[
\sin n z = 2 ^ {n - 1} \sin z \sin \left(z + \frac {\pi}{n}\right) \sin \left(z + \frac {2 \pi}{n}\right) \dots \sin \left(z + \frac {(n - 1) \pi}{n}\right).\tag{2}
\]

Indeed, one can write

\[
\begin{array}{l} \sin n z = \frac {e ^ {n i z} - e ^ {- n i z}}{2 i} = \frac {e ^ {- n i z} (e ^ {2 n i z} - 1)}{2 i} \\ \quad = \frac {e ^ {- n i z} (e ^ {2 i z} - 1) (e ^ {2 i z} - e ^ {- 2 i \pi / n}) \dots (e ^ {2 i z} - e ^ {- 2 (n - 1) i \pi / n})}{2 i} \\ \quad = A \sin z \sin \left(z + \frac {\pi}{n}\right) \sin \left(z + \frac {2 \pi}{n}\right) \dots \sin \left(z + \frac {(n - 1) \pi}{n}\right) \end{array}
\]

with

\[
\mathrm{A} = (2 i) ^ {n - 1} e ^ {- \frac {i \pi}{n} (1 + 2 + \dots + (n - 1))} = (2 i) ^ {n - 1} e ^ {- i (n - 1) \frac {\pi}{2}} = 2 ^ {n - 1}.
\]

COROLLARY 1. For every integer n one has

\[
\sin \frac {\pi}{n} \sin \frac {2 \pi}{n} \dots \sin \frac {(n - 1) \pi}{n} = \frac {n}{2 ^ {n - 1}}.\tag{3}
\]

It suffices to divide both sides of (2) by \(\sin z\) and let z tend to 0.

COROLLARY 2. For every odd integer \(n = 2m + 1\) , and every complex number z such that nz is not an integral multiple of \(\pi\) , one has

\[
\cot n z = (- 1) ^ {m} \cot z \cot \left(z + \frac {\pi}{n}\right) \dots \cot \left(z + \frac {(n - 1) \pi}{n}\right).\tag{4}
\]

Indeed, \(\sin n\left(z + \frac{\pi}{2}\right) = \sin \left(nz + \frac{\pi}{2} + m\pi\right) = (-1)^m\cos nz\) , whence, replacing \(z\) by \(z + \frac{\pi}{2}\) in (2),

\[
\cos n z = (- 1) ^ {m} 2 ^ {n - 1} \cos z \cos \left(z + \frac {\pi}{n}\right) \dots \cos \left(z + \frac {(n - 1) \pi}{n}\right)\tag{5}
\]

and the formulae (2) and (5) give (4) on division term-by-term when \(\sin nz \neq 0\) .

In all that follows we shall always assume that \(n = 2m + 1\) is an odd integer; the formula (4) can also be written

\[
\cot n z = (- 1) ^ {m} \prod_ {k = - m} ^ {m} \cot \left(z - \frac {k \pi}{n}\right).
\]

For, one has

\[
\cot \left(z - \frac {k \pi}{n}\right) = \frac {1 + \tan z \tan \frac {k \pi}{n}}{\tan z - \tan \frac {k \pi}{n}}
\]

for \(\tan z\) finite; \(\cot nz\) is thus a rational function in \(u = \tan z\) , whose numerator is of degree n - 1, and whose denominator, of degree n, has the n simple roots \(\tan k\pi/n\) ; decomposing this into simple fractions yields

\[
\cot n z = \sum_ {k = - m} ^ {m} \frac {\mathrm{A} _ {k}}{u - \tan \frac {k \pi}{n}}\tag{6}
\]

where

\[
\begin{array}{l}\mathrm{A} _ {k} = \lim _ {z \rightarrow k \pi / n} \cot n z \left(\tan z - \tan \frac {k \pi}{n}\right) = \lim _ {z \rightarrow k \pi / n} \frac {\cos n z}{\sin n z} \frac {\sin \left(z - \frac {k \pi}{n}\right)}{\cos z \cos \frac {k \pi}{n}}\\= \lim _ {h \rightarrow 0} \frac {\cos n h}{\cos \frac {k \pi}{n} \cos \left(h + \frac {k \pi}{n}\right)} \frac {\sin h}{\sin n h} = \frac {1}{n \cos^ {2} \frac {k \pi}{n}}\end{array}
\]

whence, on isolating the term in (6) corresponding to k = 0 and combining the terms corresponding to opposite values of k, and replacing z by z/n,

\[
\cot z = \frac {1}{n \tan \frac {z}{n}} + \sum_ {k = 1} ^ {m} \frac {2 n \tan \frac {z}{n}}{\cos^ {2} \frac {k \pi}{n} \left(n \tan \frac {z}{n}\right) ^ {2} - \left(n \sin \frac {k \pi}{n}\right) ^ {2}}\tag{7}
\]

valid for every complex number \(z\) not an integral multiple of \(\pi / 2\) . One can write this formula in the form

\[
\cot z = \frac {1}{n \tan \frac {z}{n}} + \sum_ {k = 1} ^ {\infty} v _ {k} (n, z)
\]

with \(v_{k}(n,z) = 0\) for \(k > m\) and

\[
v _ {k} (n, z) = \frac {2 n \tan \frac {z}{n}}{\cos^ {2} \frac {k \pi}{n} \left(n \tan \frac {z}{n}\right) ^ {2} - \left(n \sin \frac {k \pi}{n}\right) ^ {2}}
\]

for \(1 \leqslant k \leqslant m\) . We shall see that for every \(z\) contained in a compact subset K of C, not containing any integral multiple of \(\pi\) , and for every sufficiently large odd \(n\) , the series with general term \(v_{k}(n,z)\) is normally convergent. Indeed, as \(n\) tends to \(+\infty\) , \(\tan \frac{z}{n}\) tends to \(\frac{z}{n}\) uniformly on K, so there exists a number M \textgreater{} 0 such that \(\left| n \tan \frac{z}{n} \right| \leqslant M\) for every sufficiently large \(m\) and every \(z \in K\) . On the other hand, for \(0 \leqslant x \leqslant \pi/2\) one has \(\sin x/x \geqslant 1 - \frac{x^2}{6} \geqslant \frac{1}{2}\) , so for \(1 \leqslant k \leqslant m\) one has \(n \sin \frac{k\pi}{n} \geqslant k\pi/2\) ; consequently, when \(m\) is sufficiently large, for every integer \(k\) such that \(k\pi / 2 > \mathbf{M}\) one has \(|v_k(n, z)| \leqslant \frac{8\mathbf{M}}{k^2\pi^2 - 4\mathbf{M}^2}\) , which proves our assertion. For \(k\) fixed, \(v_k(n, z)\) tends (uniformly on \(\mathbf{K}\) ) to \(\frac{2z}{z^2 - k^2\pi^2}\) as \(n\) tends to \(+\infty\) . Consequently:

THEOREM 1. For every complex number z not an integral multiple of \(\pi\) one has

\[
\cot z = \frac {1}{z} + \sum_ {n = 1} ^ {\infty} \frac {2 z}{z ^ {2} - n ^ {2} \pi^ {2}}\tag{8}
\]

the series on the right being normally convergent on every compact subset \(K \subset C\) not containing any integer multiple of \(\pi\) (Eulerian expansion of \(\cot z\) ).

